\documentclass[10pt]{article}

\usepackage[margin=1in]{geometry}
\usepackage{booktabs}
\usepackage{hyperref}
\usepackage{microtype}
\usepackage{amsmath}
\usepackage{graphicx}

\title{Explaining Retrieved Evidence in Retrieval-Augmented Generation\\
with Shapley Values and Shapley Interactions}
\author{Anonymous}
\date{\today}

\newcommand{\tbd}[1]{\textbf{TBD: #1}}

\begin{document}
\maketitle

\begin{abstract}
Retrieval-augmented generation (RAG) exposes a set of retrieved passages but
does not directly explain how each passage affects the generated answer.
We formulate retrieved passages as players in a cooperative game and use
Shapley values to quantify individual evidence contributions. Shapley
interaction indices are used to identify complementary, redundant, and
conflicting evidence. The evaluation separates retrieval quality, answer
correctness, explanation faithfulness, and interaction recovery. It combines
a controlled corpus with known evidence structure, an external
scientific-question-answering benchmark, and a difficult research-paper case
study. The controlled evaluation compares lexical, target-likelihood, and
contrastive-likelihood value functions under a fixed dense retriever, with a
TF--IDF contrastive-likelihood baseline.
\end{abstract}

\section{Introduction}
RAG systems combine a parametric language model with externally retrieved
documents \cite{lewis2020rag}. Retrieval scores describe query--passage relevance, but they do not
measure how much a passage supports a particular answer. This distinction is
important when retrieved passages are distracting, mutually necessary,
duplicated, or contradictory.

This work studies post-retrieval explanation rather than proposing a new
general-purpose retriever. Given a question, a generated or reference answer,
and the retrieved top-$k$ passages, we define a cooperative game whose value
function scores an answer under subsets of the retrieved context. The resulting
Shapley values and interactions provide a chunk-level explanation.

\paragraph{Research questions.}
\begin{description}
  \item[RQ1] Do Shapley attributions identify gold supporting evidence more
  accurately than retrieval scores, leave-one-out contributions, and random
  rankings?
  \item[RQ2] Can Shapley interaction values distinguish complementary evidence
  from redundant or conflicting evidence?
\end{description}
Stability checks and answer diagnostics are retained as supporting analyses,
not standalone research questions, because the central object of study is
Shapley-based evidence attribution and interaction recovery rather than RAG
system diagnosis.

\paragraph{Contributions.}
The intended contributions are:
\begin{enumerate}
  \item a cooperative-game formulation for retrieved RAG passages;
  \item dual explanations of independently annotated gold answers and actual
  generated answers;
  \item a controlled end-to-end corpus with known direct, complementary,
  missing, conflicting, and redundant evidence patterns; and
  \item an evaluation protocol separating retrieval, generation, attribution,
  interaction, and supporting diagnostic analyses.
\end{enumerate}

\section{Background and Related Work}
\subsection{Retrieval-Augmented Generation}
Retrieval-augmented generation combines a retriever over external text with a
parametric generator \cite{lewis2020rag}. In this setting, retrieval scores
rank passages by query relevance before generation. They do not necessarily
measure whether a passage supports a specific answer, nor whether the generator
actually depends on that passage. A high retrieval score can correspond to a
distractor, a redundant paraphrase, or context that is relevant to the question
but irrelevant to the final answer. This work therefore treats retrieval
quality, answer correctness, and answer-conditioned evidence support as
separate evaluation objects.

The method is post-retrieval: it does not claim to improve dense retrieval,
chunking, or generation. Instead, it asks how the already retrieved top-$k$
chunks affect a selected value function. This is useful for RAG debugging
because the same retrieved set can be explained against two targets: an
independently annotated reference answer and the answer actually generated by
the model.

\subsection{Shapley Values}
Following the cooperative-game value introduced by Shapley
\cite{shapley1953value}, let $N$ be the set of retrieved passages and $v(S)$
the answer score when only
coalition $S \subseteq N$ is visible. The Shapley value of passage $i$ is
\begin{equation}
\phi_i(v) =
\sum_{S \subseteq N \setminus \{i\}}
\frac{|S|!(|N|-|S|-1)!}{|N|!}
\left[v(S \cup \{i\}) - v(S)\right].
\end{equation}
It averages the marginal contribution of passage $i$ over all context orders.

\subsection{Shapley Interactions}
Second-order interactions measure whether two passages have a joint effect
that is not captured by the sum of their individual effects. The implementation
uses \texttt{shapiq.KernelSHAPIQ} with \texttt{index="k-SII"} and
\texttt{max\_order=2}, which estimates the order-limited Shapley interaction
index. For a pair of passages $i$ and $j$, the relevant local effect is the
second discrete derivative
\begin{equation}
\Delta_{ij}v(S) =
v(S \cup \{i,j\}) - v(S \cup \{i\}) - v(S \cup \{j\}) + v(S),
\end{equation}
averaged over coalitions $S \subseteq N \setminus \{i,j\}$ using Shapley-style
weights. Positive values indicate that the two passages are complementary for
the selected target and value function: together they raise the score more than
their separate effects predict. Negative values indicate redundancy or
interference: adding both provides less benefit than adding them separately, or
one passage works against the contribution of the other. Following cooperative-
game interaction work \cite{grabisch1999interaction}, these signs are
interpreted relative to the chosen game; they do not certify factual truth.

Controlled interaction evaluation is only defined for labelled passage pairs
whose two members are both retrieved. If one or both labelled passages are not
in the top-$k$ retrieved set, the pair is reported as skipped rather than
counted as correct or incorrect.

\section{Method}
\subsection{RAG Pipeline}
The system parses or loads a document corpus, constructs candidate chunks,
retrieves the top-$k$ candidates, generates an answer, and then explains the
retrieved set. The controlled benchmark uses the same retrieval implementation
as the PDF path; it does not inject a preselected set of passages.

The interactive demonstration presents the retrieved context, answer targets,
and evidence effects as separate objects. Uploaded PDFs use BGE-base and top-6
retrieval by default, followed by E4B contrastive-likelihood attribution. When
a human-verified reference answer is available, the same retrieved chunks are
used to construct two independent games: one explaining support for the
reference answer and one explaining support for the generated answer. The user
can switch between their complete attribution, interaction, and faithfulness
results. Without a reference annotation, the interface reports generated-
answer attribution only. Retrieval rank is displayed separately from Shapley
rank. Bundled didactic scenarios are marked as preselected coalition
demonstrations rather than retrieval evaluations.

\subsection{Players and Coalitions}
Each retrieved chunk is one player. A coalition mask determines which chunks
are included in the model context. Chunk order and prompt formatting are held
fixed within an experiment. The controlled benchmark retrieves $k=6$ chunks,
so each game has $2^6=64$ coalitions. KernelSHAP and second-order KernelSHAPIQ
are each given the full 64-coalition budget.

\subsection{Dual Attribution Targets}
We keep three objects distinct:
\begin{itemize}
  \item the independently annotated gold answer $y^\star$;
  \item the retrieved-context model output $\hat{y}$; and
  \item the coalition value $v(S, y)$ for target $y$.
\end{itemize}
Gold-answer attribution measures support for the correct answer. Generated-
answer attribution measures context dependence of the answer actually emitted
by the model. Neither is interpreted as a certificate of factual correctness.

\subsection{Value Functions}
The implementation includes lexical answer support, target-answer likelihood,
contrastive likelihood relative to an empty context, and generated-answer
overlap. The model-backed primary analysis should use target likelihood or
contrastive likelihood; lexical support is retained as a reproducible,
model-free diagnostic.

Let $\ell(S,y)$ be the mean token log-likelihood of answer $y$ when the prompt
contains coalition $S$. With temperature $\tau=4$, target likelihood uses
\begin{equation}
v_{\mathrm{target}}(S,y)=\sigma\left(\ell(S,y)/\tau\right),
\end{equation}
whereas contrastive likelihood uses
\begin{equation}
v_{\mathrm{contrast}}(S,y)=
\sigma\left((\ell(S,y)-\ell(\emptyset,y))/\tau\right).
\end{equation}
The latter removes the answer's no-context baseline before applying the
bounded transform. The empty-context likelihood is cached within each game;
this changes runtime but not the value function.

\section{Evaluation}
\subsection{Datasets}
\paragraph{Controlled corpus.}
The bundled controlled corpus is maintained in
\path{evals/cases/controlled_benchmark.json}. It contains 108 passages and 27
questions covering five evidence patterns: direct, complementary, missing,
conflicting, and redundant evidence. Gold evidence is labelled by stable
passage IDs. The retriever searches the full 108-passage corpus for every query.
The primary dense retriever is BAAI BGE-base-en-v1.5
\cite{xiao2023cpack}, a 109.5-million-parameter English embedding model.
Queries use the model's retrieval instruction,
passages do not receive the instruction, and embeddings use the first-token
representation followed by L2 normalization.
For interaction recovery, the same controlled file contains 30 labelled
passage pairs: 10 complementary, 10 redundant, and 10 conflicting. Each
labelled pair records stable passage IDs, relation type, and expected sign.
Earlier first-order attribution runs were recorded before this interaction
expansion and are kept as historical 20-case results; new controlled runs use
the 27-case file by default.

\paragraph{External benchmark.}
The external validation uses QASPER \cite{dasigi2021qasper}, a scientific
question-answering benchmark over NLP papers released under CC BY 4.0. We use
the official validation split from \texttt{allenai/qasper} and freeze a
30-question subset in
\path{evals/cases/qasper_validation_30.json}. The frozen manifest records the
source revision
\texttt{fdc9d8214fbab5dd782958601db4d678e6934a54}, paper IDs, question IDs, and
answer annotation IDs; the full external text is downloaded into a local cache
rather than copied into the repository.

Each QASPER paper is converted into paragraph-level retrieval candidates, with
abstract paragraphs included as candidates. A case is eligible when it has an
answerable free-form or extractive annotation, a 2--40 word reference answer,
and at least one official evidence string that maps back to a paper paragraph.
The selection is deterministic: validation row order, then question order, then
annotation order, with at most one question per paper. The resulting 30 cases
contain 10 free-form and 20 extractive answers. QASPER is used as an external
domain validation, not as an interaction-label benchmark.

\paragraph{Artemis case study.}
The Artemis PDF is used as a qualitative prior-knowledge control. Its questions
target report-specific Artemis III details that the model is unlikely to answer
reliably without retrieved context, helping separate retrieval-grounded behavior
from parametric knowledge.
A lightweight frozen controlled manifest derived from the same report is stored
in \path{evals/cases/artemis_controlled_benchmark.json}; it contains 10
direct-evidence questions and 30 short passages, with no interaction labels, and
uses top-$3$ retrieval to keep model-backed attribution inexpensive. The
completed recorded run is stored in
\path{evals/runs/20260617_artemis_bge_contrastive_ll_e4b_layers20}. It is used
as qualitative supporting evidence rather than as a primary quantitative
benchmark.

\subsection{Metrics}
Retrieval is measured with Recall@$k$, Precision@$k$, and mean reciprocal rank.
Answer quality is measured separately under closed-book, gold-context, and
retrieved-context conditions. Attribution is evaluated with top-evidence
accuracy, attribution mass on gold evidence, deletion AUC, and insertion AUC.
Labelled evidence pairs are used to evaluate interaction sign recovery only in
the controlled corpus, where such pair labels exist.

\subsection{Knowledge-Source Intervention}
The same model answers each question with no context, gold context, and
retrieved context. If the model fails closed-book but succeeds with retrieved
evidence, the answer is classified as retrieval-enabled. If gold context
succeeds but retrieved context fails, the error is classified as retrieval
failure. If gold context also fails, the outcome indicates a generation or
model-capability failure.

When both closed-book and retrieved-context answers are correct, correctness
alone cannot identify whether the model used parametric knowledge or retrieved
evidence. We therefore report this case as source-not-identifiable unless
context interventions produce additional evidence of dependence.

\subsection{Experimental Setup}
The three model-backed controlled runs use the local quantized instruction model
\texttt{gemma-4-E4B-it-Q4\_K\_M.gguf} through llama.cpp with a 4096-token
context. The QASPER run uses the same model, context size, top-$6$ retrieval,
and contrastive-likelihood value function. BGE embeddings run on Apple MPS.
Twenty language-model layers are offloaded to Metal and the remaining
computation uses the CPU path. Full offload was not used because it caused a
native process failure on the 32-GB M2 Max test machine. Generation is
deterministic with temperature zero, and the knowledge-source classifier uses
token F1 $\geq 0.65$ as its operational correctness threshold.

\subsection{Baselines and Ablations}
The explanation ranking is compared with retrieval score, leave-one-out
contribution, and random order. Exact Shapley computation on small games is used
to evaluate approximation error. Chunk size, overlap, retrieval method, value
function, model, budget, and random seed are treated as robustness ablations
rather than the central contribution.

\section{Results}
\begin{table}[ht]
\centering
\small
\caption{Controlled-benchmark attribution results on the 27-case benchmark.
Hit top gold and hit gold mass are computed only on cases where at least one
labelled evidence passage was retrieved. Lower deletion AUC is better.}
\begin{tabular}{llrrrrrrr}
\toprule
Setting & Target & Cases & Gold hit & Recall@$6$ & MRR & Hit top gold &
Hit gold mass & Deletion AUC \\
\midrule
BGE-base + lexical & Gold & 27 & 23 & 0.957 & 0.967 & 0.957 & 0.697 & 0.255 \\
BGE-base + target LL & Gold & 27 & 23 & 0.957 & 0.967 & 0.957 & 0.679 & 0.038 \\
BGE-base + target LL & Generated & 27 & 23 & 0.957 & 0.967 & 0.957 & 0.812 & -- \\
BGE-base + contrastive LL & Gold & 27 & 23 & 0.957 & 0.967 & 0.957 & 0.666 & 0.046 \\
BGE-base + contrastive LL & Generated & 27 & 23 & 0.957 & 0.967 & 0.957 & 0.800 & -- \\
TF--IDF + contrastive LL & Gold & 27 & 22 & 0.891 & 0.935 & 1.000 & 0.673 & 0.041 \\
TF--IDF + contrastive LL & Generated & 27 & 22 & 0.891 & 0.935 & 0.955 & 0.780 & -- \\
\bottomrule
\end{tabular}
\end{table}

\begin{table}[ht]
\centering
\small
\caption{External QASPER validation with BGE-base retrieval, Gemma E4B, and
contrastive likelihood. Hit top gold and hit gold mass condition on at least
one labelled evidence passage appearing in the retrieved top-6 context. Lower
deletion AUC is better.}
\begin{tabular}{lrrrrrrr}
\toprule
Target & Cases & Gold hit & Recall@$6$ & MRR & Hit top gold & Hit gold mass & Deletion AUC \\
\midrule
Contrastive LL / gold & 30 & 21 & 0.550 & 0.451 & 0.857 & 0.592 & 0.033 \\
Contrastive LL / generated & 30 & 21 & 0.550 & 0.451 & 0.667 & 0.529 & -- \\
\bottomrule
\end{tabular}
\end{table}
The QASPER results in Table 2 are from
\path{evals/runs/20260615_qasper_30_bge_e4b_contrastive}.

\begin{table}[ht]
\centering
\small
\caption{Controlled attribution stability from recorded artifacts. Spearman is
computed over absolute first-order Shapley values.}
\begin{tabular}{lrrrr}
\toprule
Comparison & Target & Retrieved Jaccard & Top agreement & Spearman \\
\midrule
BGE lexical vs BGE target LL & Gold & 1.000 & 0.704 & 0.523 \\
BGE target LL vs BGE contrastive LL & Gold & 1.000 & 1.000 & 0.979 \\
BGE target LL vs BGE contrastive LL & Generated & 1.000 & 1.000 & 0.992 \\
BGE lexical vs BGE contrastive LL & Gold & 1.000 & 0.704 & 0.514 \\
BGE contrastive LL vs TF--IDF contrastive LL & Gold & 0.592 & 0.741 & 0.428 \\
BGE contrastive LL vs TF--IDF contrastive LL & Generated & 0.592 & 0.778 & 0.467 \\
\bottomrule
\end{tabular}
\end{table}
The stability results in Table 3 are from
\path{evals/runs/20260616_controlled_stability}.

\subsection{RQ1: Evidence Identification}
The 27-case BGE-base lexical run,
\path{evals/runs/20260616_controlled_interactions_bge_lexical}, retrieved at
least one labelled evidence passage in 23 cases. Evidence Recall@6 was 0.957
and MRR was 0.967. Conditional on those gold-hit cases, the highest absolute
Shapley attribution was assigned to labelled evidence in 22 of 23 cases
(0.957), and the mean absolute attribution mass on labelled evidence was 0.697.
Shapley-order deletion AUC was 0.255.

The completed E4B target-likelihood run,
\path{evals/runs/20260616_controlled_bge_target_ll_e4b_layers20}, and the E4B
contrastive-likelihood run,
\path{evals/runs/20260616_controlled_bge_contrastive_ll_e4b_layers20}, use the
same BGE retrieved sets. For the gold-answer target, both placed labelled
evidence first in 22 of 23 gold-hit cases (0.957), with hit-conditioned
labelled-evidence masses of 0.679 and 0.666. Their deletion AUCs were 0.038 and
0.046. For the generated-answer target, both again placed labelled evidence
first in 22 of 23 gold-hit cases, while assigning larger labelled-evidence
masses of 0.812 and 0.800. On the controlled corpus, generated-answer behavior
is therefore still strongly aligned with the labelled evidence.

Holding E4B and the contrastive value function fixed, TF--IDF retrieved less
labelled evidence than BGE on the current benchmark:
\path{evals/runs/20260616_controlled_tfidf_contrastive_ll_e4b_layers20}
retrieved at least one labelled evidence passage in 22 rather than 23 cases,
with Recall@6 of 0.891 and MRR of 0.935. Conditional on the gold-hit cases,
TF--IDF nevertheless placed labelled evidence first for the gold-answer target
in all 22 cases (1.000), with hit-conditioned labelled-evidence mass of 0.673.
For the generated-answer target, top-gold accuracy was 21/22 (0.955) and
labelled-evidence mass was 0.780. Thus the retriever ablation mainly reduces
evidence coverage; once labelled evidence is retrieved, first-order attribution
remains comparable on this synthetic benchmark.

On the 30-case QASPER validation subset, the same BGE + E4B contrastive
configuration achieved Recall@6 of 0.550 and MRR of 0.451. Because Shapley
cannot assign attribution to evidence that was not retrieved, the most
informative external result conditions on the 21 cases where at least one
labelled evidence passage appears in the retrieved context. In this gold-hit
subset, the highest gold-answer Shapley attribution was assigned to labelled
evidence in 18 of 21 cases (0.857), with mean gold attribution mass of 0.592.
Generated-answer attribution was less aligned with the external evidence
labels: the top generated-answer attribution was labelled evidence in 14 of
21 gold-hit cases (0.667), with mean labelled-evidence mass of 0.529. These
results show that the controlled corpus does not by itself estimate
performance on scientific papers, but also that first-order Shapley attribution
often recovers labelled evidence when retrieval has surfaced it.

\subsection{RQ2: Interaction Recovery}
Table~\ref{tab:controlled-interactions} reports controlled interaction
recovery. Interaction coverage is the fraction of labelled pairs for which
both passages were retrieved in the top-6 context. Sign recovery is conditional
on coverage: skipped pairs are not counted as sign failures because the
interaction game is undefined for passages that were not retrieved.

\begin{table}[ht]
\centering
\small
\caption{Controlled interaction recovery. Sign recovery is conditional on both
passages being retrieved.}
\label{tab:controlled-interactions}
\resizebox{\textwidth}{!}{%
\begin{tabular}{lr|rrrr|rrrr|rrrr}
\toprule
Type & Labelled
& \multicolumn{4}{c|}{Lexical / Gold}
& \multicolumn{4}{c|}{Contrastive LL / Gold}
& \multicolumn{4}{c}{Contrastive LL / Generated} \\
\cmidrule(lr){3-6}\cmidrule(lr){7-10}\cmidrule(lr){11-14}
& & Eval. & Cov. & Correct & Recovery
& Eval. & Cov. & Correct & Recovery
& Eval. & Cov. & Correct & Recovery \\
\midrule
Complementary & 10 & 9 & 0.900 & 5 & 0.556 & 9 & 0.900 & 3 & 0.333 & 9 & 0.900 & 5 & 0.556 \\
Redundant & 10 & 9 & 0.900 & 9 & 1.000 & 9 & 0.900 & 8 & 0.889 & 9 & 0.900 & 3 & 0.333 \\
Conflicting & 10 & 10 & 1.000 & 10 & 1.000 & 10 & 1.000 & 5 & 0.500 & 10 & 1.000 & 1 & 0.100 \\
\bottomrule
\end{tabular}
}
\end{table}

The lexical run is
\path{evals/runs/20260616_controlled_interactions_bge_lexical}. The
model-backed contrastive-likelihood run is
\path{evals/runs/20260616_controlled_bge_contrastive_ll_e4b_layers20}. The
derived RQ2 table is
\path{evals/runs/20260616_controlled_interaction_summary}. The
lexical value function recovers redundant and conflicting negative signs well
on this controlled set, but complementary signs remain harder: only five of
nine retrieved complementary pairs have the expected positive sign. This
supports using lexical support as a transparent diagnostic baseline, not as the
strongest evidence for interaction recovery. E4B contrastive likelihood with
20 Metal-offloaded layers produces strong first-order attribution metrics, but
weaker pairwise sign recovery on the expanded interaction labels: 16 of 28
evaluable gold-target signs and 9 of 28 generated-target signs match the
expected relation signs. This suggests that the current sign convention and
value function capture some redundancy but do not reliably separate all
complementary and conflicting cases. A full Metal offload attempt
(\texttt{n\_gpu\_layers=-1}) failed with a native llama.cpp process crash, so
the completed model-backed RQ2 run uses partial Metal offload
(\texttt{n\_gpu\_layers=20}).

\subsection{Supporting Analysis: Attribution Stability}
We analyze stability from the completed controlled-run artifacts without
re-running model scoring. Chunks are ranked by absolute first-order Shapley
attribution. With the retrieved set fixed, target likelihood and contrastive
likelihood produce nearly identical rankings: Spearman correlation is 0.979 for
gold-answer attribution and 0.992 for generated-answer attribution, with
top-attribution agreement of 1.000 in both cases. Thus the main likelihood-
based attribution conclusion is stable to subtracting the empty-context
baseline.

Lexical support is less stable relative to contrastive likelihood even with the
same BGE retrieved chunks. Its gold-answer Spearman correlation with
contrastive likelihood is 0.514 and top-attribution agreement is 0.704. This
supports treating lexical support as a transparent diagnostic baseline rather
than the primary value function.

Changing the retriever is the largest observed source of instability. Comparing
BGE + contrastive likelihood with TF--IDF + contrastive likelihood, the mean
retrieved-set Jaccard overlap is 0.592. On the union of retrieved chunks, the
gold-answer Spearman correlation drops to 0.428 and generated-answer
correlation to 0.467. Therefore explanation stability depends strongly on the
upstream retrieval set, not only on the Shapley estimator or value function.

\subsection{Supporting Analysis: Answer Diagnostics}
Generation outputs are identical across the two BGE likelihood value
functions because they share the same retriever, prompt, and E4B generator.
For BGE contrastive likelihood on the 27-case controlled benchmark, eight cases
were classified as retrieval-enabled, one as retrieval failure, 14 as
generation-or-model failure, and four as missing evidence by construction. Mean
token F1 was 0.008 closed-book, 0.503 with annotated gold context, and 0.560
with retrieved context.

Replacing BGE with TF--IDF changed the retrieved-context mean F1 to 0.596 and
produced the same coarse source classification on this small controlled set.
Closed-book and gold-context results were unchanged, so the retrieved-context
difference is attributable to the changed retrieved set under this operational
protocol.
These labels use the 0.65 token-F1 threshold and are not direct observations
of the model's internal knowledge source.

For QASPER, mean token F1 was 0.026 closed-book, 0.292 with annotated gold
context, and 0.178 with retrieved context. These token-F1 values are retained
as answer-quality diagnostics, not as the primary Shapley result. The 0.65
threshold marks 29 cases as gold-context low-F1 flags and one as a
retrieved-context low-F1 flag, but short extractive references make this
threshold sensitive to over-answering. The attribution analysis in RQ1 is
therefore interpreted primarily through retrieval coverage, top-gold
attribution, and labelled-evidence attribution mass rather than through the
F1-derived diagnostic labels.

The Artemis PDF case study further probes parametric-knowledge confounding.
In the 10 recorded Artemis artifacts from
\path{evals/runs/20260617_artemis_bge_contrastive_ll_e4b_layers20}, mean
closed-book, gold-context, and retrieved-context F1 were 0.019, 0.760, and
0.769. Retrieval found labelled evidence in all cases (Recall@3 = 1.000), the
generated-answer target ranked a labelled passage first in 10/10 cases, and the
mean generated-answer labelled-evidence attribution mass was 0.889. For
\texttt{artemis\_first\_surface\_mission}, closed-book F1 was 0.000,
retrieved-context F1 was 0.872, and the generated-answer Shapley attribution
assigned 0.096 to the retrieved gold passage, corresponding to 0.685 of the
positive generated-attribution mass. This is the intended diagnostic role of the
generated-answer target: it checks whether the answer the model actually
produced is supported by retrieved evidence, while the gold-answer target asks a
separate reference-support question.

\section{Discussion}
The explanation is conditional on the chosen answer target, retrieved set,
prompt, model, and value function. A high Shapley value means that a passage
changes the selected score; it does not prove that the passage is true.
Likewise, deletion and insertion evaluate faithfulness to the value function,
not factual correctness.

\section{Limitations}
\begin{itemize}
  \item The controlled corpus is synthetic by construction and may overstate
  performance relative to naturally occurring scientific documents.
  \item Evidence annotations are not unique when several passages express the
  same fact.
  \item The QASPER subset uses one deterministic annotation per question and
  paragraph-level evidence mapping; alternative annotator aggregation or
  evidence-span chunking could change the external metrics.
  \item Model knowledge and retrieved knowledge are not fully identifiable
  when both closed-book and context-conditioned answers are correct.
  \item Coalition evaluation grows exponentially with the retrieved set, so
  larger games require approximation.
  \item PDF parsing, OCR, tables, equations, and chunk boundaries can introduce
  retrieval errors unrelated to the attribution method.
  \item Likelihood-based value functions depend on model calibration and exact
  answer wording.
  \item The local E4B model used partial rather than full Metal offload because
  full offload exceeded the stable native runtime configuration on the test
  machine; runtime configuration may affect floating-point results.
\end{itemize}

\section{Conclusion}
This project evaluates Shapley values as explanations of retrieved evidence,
not as a claim that a RAG system is perfectly accurate. The proposed protocol
separates the quality of retrieval, answers, value functions, and explanations,
making both successful and failed RAG runs scientifically informative.

\section*{Research Log}
\begin{description}
  \item[2026-06-17, research-question framing] Refocused the manuscript and
  HTML report around two core Shapley research questions: first-order evidence
  identification and pairwise interaction recovery. Stability checks and
  answer-source diagnostics are retained as supporting analyses rather than
  standalone research questions, because they characterize the experimental
  setup but are not the central contribution.
  \item[2026-06-13] Defined four research questions and separated retrieval,
  answer, attribution, and interaction evaluation. Added a 20-question,
  80-passage controlled corpus searched by the real retriever. Defined dual
  gold-answer and generated-answer attribution targets and the closed-book,
  gold-context, retrieved-context knowledge-source protocol. Numerical results
  are pending the first recorded run.
  \item[2026-06-14, experiment reset] Removed the two exploratory TF--IDF run
  outputs and their numerical claims. Fixed the controlled experiment matrix
  to BGE-base with lexical, target-likelihood, and contrastive-likelihood value
  functions, plus TF--IDF with contrastive likelihood. The first recorded run
  is BGE-base + lexical.
  \item[2026-06-14, BGE-base + lexical] Completed all 20 controlled cases in
  \path{evals/runs/20260614_bge_base_lexical}. BGE-base-en-v1.5 achieved
  evidence Recall@6 of 0.906 and MRR of 0.953. Lexical Shapley attribution
  achieved top-gold accuracy of 0.938, gold attribution mass of 0.629, and
  deletion AUC of 0.228, versus 0.244 for retrieval-score order and 0.460 for
  random order. Interaction sign accuracy was 0.600 across five evaluable
  pairs, with weaker recovery for complementary than redundant evidence.
  \item[2026-06-14, E4B likelihood runs] Completed 20 controlled cases for
  BGE-base + target likelihood in
  \path{evals/runs/20260614_bge_base_target_ll_e4b} and BGE-base + contrastive
  likelihood in
  \path{evals/runs/20260614_bge_base_contrastive_ll_e4b}. Both used the local
  E4B Q4 model with 20 layers offloaded to Metal; full offload caused a native
  process failure on the 32-GB M2 Max. Target and contrastive likelihood
  achieved top-gold accuracy of 1.000, gold masses of 0.669 and 0.655,
  deletion AUCs of 0.035 and 0.038, and interaction sign accuracy of 1.000
  across five evaluable pairs. Generated-answer top-gold accuracy was 0.938
  for both. The shared generation protocol classified eight cases as
  retrieval-enabled, one as retrieval failure, seven as generation-or-model
  failure, and four as missing evidence.
  \item[2026-06-14, TF--IDF + E4B contrastive likelihood] Completed all 20
  controlled cases in
  \path{evals/runs/20260614_tfidf_contrastive_ll_e4b}. Relative to BGE with the
  same contrastive value function and E4B model, TF--IDF reduced Recall@6 from
  0.906 to 0.844, MRR from 0.953 to 0.906, top-gold accuracy from 1.000 to
  0.938, and gold attribution mass from 0.655 to 0.613. Generated-answer
  top-gold accuracy was 0.875 and generated gold mass was 0.732. Interaction
  sign accuracy was 0.750 over four evaluable pairs. The knowledge-source
  protocol classified seven cases as retrieval-enabled, two as retrieval
  failures, seven as generation-or-model failures, and four as missing
  evidence.
  \item[2026-06-14, research-demo interface] Aligned the interactive dashboard
  with the recorded experiment protocol. The default PDF configuration is now
  BGE-base top-6 retrieval, Gemma E4B with 20 Metal-offloaded layers, and
  contrastive likelihood. The result view separates retrieval, answer-target
  selection, and Shapley attribution; distinguishes reference-answer from
  generated-answer attribution; and shows retrieval rank separately from
  Shapley rank. Preselected sample scenarios are explicitly labelled as
  coalition demonstrations rather than retrieval evaluations.
  \item[2026-06-14, dual-target interface] Extended each interactive run to
  carry independent reference-answer and generated-answer games over the same
  retrieved chunks. Controlled samples compute and display both targets. PDF
  runs expose an optional human-verified reference-answer field; when supplied,
  the interface switches the full attribution ranking, interaction matrix, and
  faithfulness audit between gold/reference support and generated-answer
  support. The header subtitle and numbered process labels were removed to
  reduce presentation ambiguity.
  \item[2026-06-15--16, QASPER external validation] Added QASPER validation
  ingestion from \texttt{allenai/qasper}, froze a 30-case deterministic subset
  in \path{evals/cases/qasper_validation_30.json}, and completed a five-case
  smoke run in
  \path{evals/runs/20260615_qasper_smoke_5_bge_e4b_contrastive}. The full
  30-case BGE-base + E4B contrastive-likelihood run is
  \path{evals/runs/20260615_qasper_30_bge_e4b_contrastive}. It achieved
  Recall@6 of 0.550, MRR of 0.451, gold-answer top-gold attribution of 0.600,
  gold attribution mass of 0.415, generated-answer top-gold attribution of
  0.467, and generated gold mass of 0.370. In the retrieval-conditioned subset
  where at least one labelled evidence passage was retrieved (21/30 cases),
  gold-answer top-gold attribution was 18/21 (0.857) with mean labelled-evidence
  mass of 0.592; generated-answer top-gold attribution was 14/21 (0.667) with
  mean labelled-evidence mass of 0.529. Mean closed-book, gold-context, and
  retrieved-context token F1 were 0.026, 0.292, and 0.178, respectively; these
  token-F1 values are retained as answer diagnostics rather than the primary
  Shapley result.
  \item[2026-06-16, controlled stability analysis] Added an artifact-only
  stability analysis in \path{evals/runs/20260616_controlled_stability}, later
  regenerated from the completed 27-case controlled runs. With BGE retrieval
  fixed, target likelihood and contrastive likelihood were nearly identical:
  top-attribution agreement was 1.000 and Spearman correlation was 0.979 for
  gold-answer attribution and 0.992 for generated-answer attribution. Lexical
  versus contrastive likelihood had gold-answer Spearman correlation of 0.514
  and top-attribution agreement of 0.704. Changing the retriever from BGE to
  TF--IDF reduced retrieved-set Jaccard overlap to 0.592 and union-ranked
  Spearman correlation to 0.428 for gold-answer attribution and 0.467 for
  generated-answer attribution.
  \item[2026-06-16, interaction-aware controlled benchmark] Updated
  \path{evals/cases/controlled_benchmark.json} from the original 20-case,
  80-passage corpus to the current 27-case, 108-passage controlled benchmark.
  The file now includes 30 labelled interaction pairs: 10 complementary, 10
  redundant, and 10 conflicting. Added artifact-level reporting of interaction
  coverage, skipped
  pairs, conditional sign recovery, full pairwise interactions, coalition
  scores, and separate generated-answer interaction diagnostics when a model
  backend is available. The completed BGE-base + lexical interaction run is
  \path{evals/runs/20260616_controlled_interactions_bge_lexical}; it
  retrieved both passages for 28 of 30 labelled pairs and recovered 24 of 28
  signs. Per type, recovery was 5/9 for complementary, 9/9 for redundant, and
  10/10 for conflicting pairs. The derived RQ2 table is
  \path{evals/runs/20260616_controlled_interaction_summary}.
  Initial E4B likelihood attempts showed that full Metal offload crashed in
  llama.cpp, and a CPU-only one-case smoke run was interrupted after more than
  two minutes without completing the first case.
  \item[2026-06-17, Artemis prior-knowledge control] Added
  \path{evals/cases/artemis_controlled_benchmark.json}, a lightweight
  10-question, 30-passage direct-evidence manifest derived from the Artemis III
  Science Definition Team PDF. The file intentionally contains no interaction
  labels, uses top-3 retrieval to keep model-backed attribution inexpensive, and
  is intended as a qualitative prior-knowledge control rather than a primary
  quantitative benchmark: the questions target report-specific details that
  should be difficult to answer reliably without retrieved context. The
  completed BGE-base + Gemma E4B contrastive-likelihood run is
  \path{evals/runs/20260617_artemis_bge_contrastive_ll_e4b_layers20}. It
  achieved retrieval Recall@k 1.000, MRR 0.950, gold-answer top-gold
  attribution 0.900 with mean labelled-evidence mass 0.834, and
  generated-answer top-gold attribution 1.000 with mean labelled-evidence mass
  0.889. Mean closed-book, gold-context, and retrieved-context token F1 were
  0.019, 0.760, and 0.769, respectively.
  \item[2026-06-16, E4B contrastive interaction completion] Completed the
  27-case controlled BGE-base + E4B contrastive-likelihood run with BGE
  embeddings on MPS and 20 Gemma E4B layers offloaded through llama.cpp Metal:
  \path{evals/runs/20260616_controlled_bge_contrastive_ll_e4b_layers20}. The
  run was interrupted after 18 case artifacts had been written, then resumed by
  running the nine missing case IDs and rebuilding \path{summary.csv} from all
  27 artifacts. The rebuilt manifest is marked
  \texttt{rebuilt\_from\_artifacts}. The derived interaction table in
  \path{evals/runs/20260616_controlled_interaction_summary} reports 16/28
  correct gold-target signs and 9/28 correct generated-target signs. Full
  Metal offload with \texttt{n\_gpu\_layers=-1} still caused native llama.cpp
  crashes, so the completed run uses partial Metal offload.
  \item[2026-06-17, completed missing 27-case attribution ablations] Added the
  two missing controlled-benchmark runs on the current 27-case corpus:
  BGE-base + E4B target likelihood in
  \path{evals/runs/20260616_controlled_bge_target_ll_e4b_layers20} and
  TF--IDF + E4B contrastive likelihood in
  \path{evals/runs/20260616_controlled_tfidf_contrastive_ll_e4b_layers20}.
  The BGE target-likelihood run retrieved labelled evidence in 23/27 cases,
  with Recall@6 of 0.957, MRR of 0.967, hit-conditioned gold-answer top-gold
  attribution of 0.957, hit-conditioned gold mass of 0.679, and deletion AUC
  of 0.038. Its generated-answer hit-conditioned top-gold attribution was
  0.957 with gold mass 0.812. The TF--IDF contrastive run retrieved labelled
  evidence in 22/27 cases, with Recall@6 of 0.891, MRR of 0.935,
  hit-conditioned gold-answer top-gold attribution of 1.000, hit-conditioned
  gold mass of 0.673, and deletion AUC of 0.041. Its generated-answer
  hit-conditioned top-gold attribution was 0.955 with gold mass 0.780. The HTML
  report and manuscript tables were updated to use these artifact-derived
  values. The stability analysis was regenerated from the four current
  27-case controlled runs in
  \path{evals/runs/20260616_controlled_stability}.
\end{description}


\bibliographystyle{plain}
\bibliography{references}
\end{document}
